\subsubsection{Cálculo estadístico elemental}
\label{sec:pie-rey-estadistica}

Sea $X_{kj}$ la lectura de la posición $j$ en el punto de calibración $k$, con
$j = 1, 2, 3$ correspondientes a las posiciones superior, central e inferior. Antes
de introducir cualquier corrección se calcula la estadística directa de esas
mediciones.

El promedio de las $n = 3$ posiciones de un punto se obtiene aplicando la media
muestral de la ecuación~\eqref{eq:media-muestral}, y su dispersión con la varianza
insesgada de la ecuación~\eqref{eq:varianza-muestral-insesgada}. La incertidumbre de
ese promedio sigue la desviación estándar de la media deducida en la
ecuación~\eqref{eq:sigma-media-experimental}.

Lo que no tiene precedente en el desarrollo anterior es el error de indicación
elemental, que se define como la diferencia entre el promedio y el valor nominal del
patrón, sin corrección térmica ni de certificado:

\begin{equation}
    \therefore\quad
    E_k^{\text{simple}} = \bk{X}_n - VP_k .
    \label{eq:E-simple}
\end{equation}

La Tabla~\ref{tab:pie-rey-estadistica} recoge estos resultados para los cinco puntos.
Los errores simples alcanzan $33.3\ \mu\text{m}$ en los puntos de $150$ mm y $200$ mm
y se anulan en los extremos del alcance, donde las tres lecturas coinciden.

\begin{table}[H]
\centering
\caption{Cálculo estadístico elemental por punto de calibración. $\bk{X}_n$ es la media de las tres posiciones, $s$ la desviación estándar muestral y $E^{\text{simple}}$ el error de indicación sin correcciones.}
\label{tab:pie-rey-estadistica}
\begin{tabular}{rrrrr}
\toprule
$VP$ (mm) & $\bk{X}_n$ (mm) & $s$ (mm) & $s/\sqrt{3}$ (mm) & $E^{\text{simple}}$ ($\mu$m) \\
\midrule
50 & 50.0000 & 0.0000 & 0.0000 & 0.0 \\
100 & 100.0167 & 0.0289 & 0.0167 & 16.7 \\
150 & 150.0333 & 0.0289 & 0.0167 & 33.3 \\
200 & 200.0333 & 0.0289 & 0.0167 & 33.3 \\
250 & 250.0000 & 0.0000 & 0.0000 & 0.0 \\
\bottomrule
\end{tabular}
\end{table}


\subsubsection*{Dispersión nula e insuficiencia del cálculo por punto}

Las desviaciones estándar resultan nulas en $50$ mm y en $250$ mm, y de
$0.0289$ mm en los tres puntos restantes. Un valor nulo no indica ausencia de
dispersión: indica que las tres lecturas coinciden exactamente dentro de la
graduación del instrumento.

La resolución del pie de rey es $d = 0.05$ mm, ecuación~\eqref{eq:resolucion-vernier},
de modo que dos lecturas de una misma posición solo pueden diferir en múltiplos
exactos de esa cantidad y la diferencia entre la mayor y la menor está acotada por
el paso de graduación. Con tres repeticiones, la varianza muestral de la
ecuación~\eqref{eq:varianza-muestral-insesgada}, evaluada sobre diferencias que no
pueden exceder $d$ en valor absoluto, solo admite dos valores extremos:

\begin{equation}
    s \in \left\{ 0,\ \frac{d}{\sqrt{3}} \right\}
    \quad\Longrightarrow\quad
    s \in \{ 0,\ 0.0289 \}\ \text{mm},
    \label{eq:s-posibles}
\end{equation}

que es lo que se observa en la tabla. La estadística tipo~A no puede, por tanto,
caracterizar la repetibilidad del instrumento con tres lecturas por punto: la
resolución es demasiado gruesa frente a la dispersión del pie de rey.

\subsubsection*{Repetibilidad del instrumento}

Sea $R_p$ la lectura repetida $p$-ésima en un mismo punto, con $n = 5$ repeticiones
consecutivas en $250$ mm. El rango de las repeticiones es

\begin{equation}
    R = \max_p R_p - \min_p R_p,
    \label{eq:rango-repetidas}
\end{equation}

y la dispersión del conjunto se calcula con la varianza insesgada de la
ecuación~\eqref{eq:varianza-muestral-insesgada}, mientras que la incertidumbre de su
promedio sigue la ecuación~\eqref{eq:sigma-media-experimental},

\begin{equation}
    \therefore\quad
    \sigma(L_{i_R}) = \frac{s_R}{\sqrt{n}} .
    \label{eq:sigma-LiR-simple}
\end{equation}

Con las cinco lecturas de $250$ mm se obtiene $R = 50\ \mu\text{m}$ y
$s_R = 27.39\ \mu\text{m}$, de modo que

\begin{equation}
    \sigma(L_{i_R}) = \frac{27.39}{\sqrt{5}}\ \mu\text{m}
    \approx 12.25\ \mu\text{m}.
    \label{eq:sigma-LiR-simple-valor}
\end{equation}

Esta es la única estimación tipo~A aprovechable del conjunto.

\subsubsection*{Alcance y límites del cálculo elemental}

El cálculo anterior no distingue el origen de las diferencias observadas. El error de
indicación de la ecuación~\eqref{eq:E-simple} mezcla el comportamiento del
instrumento con el estado del patrón y con las condiciones ambientales en que se
realizaron las mediciones: un valor no nulo puede deberse al pie de rey, a la pila de
bloques o a la dilatación térmica de ambos cuerpos. Tampoco permite asignar una
incertidumbre al resultado, porque la dispersión por punto es incompatible con la
resolución del instrumento, como muestra la ecuación~\eqref{eq:s-posibles}.

El modelo que se desarrolla a continuación descompone el error en sus contribuciones
—de contacto, de repetibilidad, de escala, geométricas, de resolución y térmicas— y
permite propagar la incertidumbre de cada una hasta el resultado final.
